\documentclass[12pt,pstricks,border=4pt]{standalone}
\usepackage{amsmath}
\usepackage[dvipsnames]{xcolor}
\usepackage{pst-optexp}

\definecolor{Probe}{HTML}{D99A00}
\definecolor{ProbePale}{HTML}{F7E8BB}
\definecolor{Reference}{HTML}{666D73}
\definecolor{ReferencePale}{HTML}{E7E9EA}
\definecolor{Common}{HTML}{2878B5}
\definecolor{CommonPale}{HTML}{B8D5EA}
\definecolor{Faraday}{HTML}{009E73}
\definecolor{FaradayPale}{HTML}{8FD3BE}
\definecolor{Affected}{HTML}{497D91}
\definecolor{AffectedPale}{HTML}{B8D7DD}
\definecolor{Mixed}{HTML}{4C7F82}
\definecolor{MixedPale}{HTML}{C5DEDC}
\definecolor{PortH}{HTML}{C03C8B}
\definecolor{PortHPale}{HTML}{B7DADB}
\definecolor{PortV}{HTML}{5146A3}
\definecolor{PortVPale}{HTML}{D3D3E0}
\definecolor{Glass}{HTML}{D9EEF3}
\definecolor{Dot}{HTML}{76538B}
\definecolor{DotPale}{HTML}{D9CBE2}
\definecolor{Ink}{HTML}{202428}
\definecolor{SoftInk}{HTML}{60666B}

\psset{linecolor=Ink,linewidth=0.8pt}
\newpsstyle{OptComp}{linecolor=Ink,linewidth=0.9pt,fillstyle=solid,fillcolor=Glass}

\def\AtomsGlyph{%
  \psellipse[fillstyle=solid,fillcolor=black!18,linecolor=Ink,linewidth=0.8pt](0,0)(0.14,0.31)%
}

\def\PhaseDotGlyph{%
  \psline[linecolor=SoftInk,linestyle=dotted,dotsep=1.2pt,linewidth=0.6pt](0,-1.05)(0,1.05)%
  \pscircle[fillstyle=solid,fillcolor=Dot,linecolor=Ink,linewidth=0.7pt](0,0){0.125}%
}

\def\StopGlyph{%
  \psline[linecolor=SoftInk,linestyle=dotted,dotsep=1.2pt,linewidth=0.6pt](0,-1.05)(0,1.05)%
  \pscircle[fillstyle=solid,fillcolor=Ink,linecolor=Ink](0,0){0.12}%
}

\def\BrightCameraGlyph{%
  \psframe[fillstyle=solid,fillcolor=black!20,linecolor=Ink,linewidth=0.9pt](0,0)(1.65,1.55)%
  \psellipse[fillstyle=solid,fillcolor=white,linestyle=none](0.82,0.78)(0.22,0.51)%
}

\def\DarkCameraGlyph{%
  \psframe[fillstyle=solid,fillcolor=Ink,linecolor=Ink,linewidth=0.9pt](0,0)(1.65,1.55)%
  \psellipse[fillstyle=solid,fillcolor=black!22,linestyle=none](0.82,0.78)(0.22,0.51)%
}

\def\FaradayBeforeGlyph{%
  \pscircle[fillstyle=solid,fillcolor=ProbePale,linecolor=Ink,linewidth=0.6pt](0,0){0.49}%
  \psline[linecolor=Probe,linewidth=1.5pt]{<->}(-0.27,0)(0.27,0)%
}

\def\FaradayAfterGlyph{%
  \pscircle[fillstyle=solid,fillcolor=ReferencePale,linecolor=Ink,linewidth=0.6pt](0,0){0.49}%
  \psline[linecolor=Common,linewidth=0.9pt]{<->}(-0.27,0)(0.27,0)%
  \rput{24}{\psline[linecolor=Faraday,linewidth=1.5pt]{<->}(-0.27,0)(0.27,0)}%
}

\def\PBSGlyph{%
  \psframe[fillstyle=solid,fillcolor=Glass,linecolor=Ink,linewidth=0.9pt](-0.39,-0.39)(0.39,0.39)%
  \psline[linecolor=Ink,linewidth=0.8pt](-0.39,0.39)(0.39,-0.39)%
}

\def\MirrorGlyph{%
  \psline[linecolor=Ink,linewidth=2.3pt](-0.50,0.50)(0.50,-0.50)%
}


\begin{document}
\psset{unit=1cm}
\begin{pspicture}(0.15,0.68)(7.42,3.38)
  \rput[l](0.40,3.10){\textbf{\small PCI}}

  \pnode(1.10,2.15){Obj}
  \pnode(7.10,2.15){Cam}

  \begin{optexp}
    \optdipole[compname=Atoms,optdipolesize=0.04 1.20,position=start]
      (Obj)(Cam){\AtomsGlyph}{}
    \lens[compname=Lone,abspos=1.50,lensheight=1.75,lensradius=2.00]
      (Obj)(Cam){}
    \optdipole[compname=Mask,optdipolesize=0.02 1.82,abspos=3.00]
      (Obj)(Cam){\PhaseDotGlyph}{}
    \lens[compname=Ltwo,abspos=4.50,lensheight=1.75,lensradius=2.00]
      (Obj)(Cam){}

    \backlayer{%
      % Yellow is reserved for the incident probe before the atoms.
      \pspolygon[fillstyle=solid,fillcolor=ProbePale,linestyle=none]
        (0.40,1.65)(1.10,1.65)(1.10,2.65)(0.40,2.65)%
      % The unaffected probe keeps the same yellow fill up to the phase dot.
      \pspolygon[fillstyle=solid,fillcolor=ProbePale,linestyle=none]
        (1.10,1.65)(2.60,1.65)(4.10,2.15)(2.60,2.65)(1.10,2.65)%
      % Purple identifies the same reference amplitude after the phase shift;
      % the fill lightness is not used as an intensity scale.
      \pspolygon[fillstyle=solid,fillcolor=DotPale,linestyle=none]
        (4.10,2.15)(5.60,1.65)(7.10,1.65)(7.10,2.65)(5.60,2.65)%
      % PCI does not resolve the two output polarisations.  One soft cyan layer
      % therefore represents the complete atom-modified transmitted field.
      \pspolygon[fillstyle=solid,fillcolor=AffectedPale,opacity=0.60,linestyle=none]
        (1.10,1.93)(1.10,2.37)(2.60,3.03)(5.60,3.03)%
        (7.10,2.37)(7.10,1.93)(5.60,1.27)(2.60,1.27)%
    }
  \end{optexp}

  \psline[linecolor=Ink,linewidth=1.2pt](7.10,1.35)(7.10,2.95)

  \rput[t](4.10,0.95){\small phase dot}

\end{pspicture}
\end{document}
